\documentclass[11pt]{article}
\usepackage{amsmath,amssymb,graphicx,hyperref,geometry}
\geometry{margin=1in}
\title{Replication Report: Split-Peak Signature of a PDW-Driven Period-8 CDW\\
in the Cuprate Vortex Halo\\ \large (Dai, Zhang, Senthil \& Lee, arXiv:1802.03009v2)}
\author{TEXTURES-100 automated replication (class: loop-current)}
\date{\today}
\begin{document}
\maketitle

\begin{abstract}
We replicate the central experimentally-testable prediction of Dai, Zhang,
Senthil \& Lee (arXiv:1802.03009v2): that a pair-density-wave (PDW)-driven
period-8 charge-density-wave (CDW) in the halo of a cuprate superconducting
vortex inherits the $2\pi$ phase winding of the $d$-wave order parameter,
producing (i) a real-space \emph{line of zeros} through the vortex core,
(ii) a \emph{split double peak} in the Fourier transform at $q\simeq Q_a/2$,
and (iii) a sign change of $\mathrm{Re}\,\tilde\nu(q)$ across $q=Q_a/2$; whereas
the CDW-driven scenario gives a single, unsplit peak. Building the paper's
real-space order-parameter fields on a coarse $256^2$ grid and Fourier
transforming them, we reproduce all four qualitative signatures
(\textbf{4/4 checks passed}).
\end{abstract}

\section{Class-mislabel disclosure}
This paper is filed under the TEXTURES-100 \emph{loop-current} class, but it is
\textbf{not} a kagome loop-current / orbital-current order theory. Orbital loop
currents appear only as a one-line citation in the introduction (intra-cell
magnetic moments detected by neutron scattering, interpreted via orbital loop
currents; refs.~6--7 in the paper). The physics of the paper is cuprate
PDW/CDW phenomenology near a superconducting vortex. The assigned kernel,
\texttt{loop\_current\_meanfield\_kernel.py} (Ollie), is a kagome Peierls-flux
tight-binding probe --- the wrong model class. Following the replication-execution
mislabel guard, we replicate the paper's \emph{actual} minimal model and headline
prediction, and credit the Ollie kernel only for methodological provenance
(real-space one-body order-parameter field $\to$ momentum-space signature).

\section{Model and method}
Near an isolated $d$-wave vortex we take the paper's profiles (Sec.~IV):
\begin{align}
\Delta_D(r,\theta) &= |\Delta_D|\,\frac{r}{\sqrt{r^2+r_0^2}}\,e^{i\theta},
   \qquad r_0 = 3.5a,\\
\Delta_P(r) &\sim \exp\!\big(1-\sqrt{r^2+\xi^2}/\xi\big),\qquad \xi = 15a.
\end{align}
The PDW-driven induced $Q_a/2$ (period-8) CDW amplitude (paper Eq.~9/14) is
\begin{equation}
\rho_\alpha(r) = F(r)\,\cos(\theta-\theta_a)\,\cos(Q_\alpha\!\cdot r + \phi_\alpha),
\qquad F(r)\sim 2c\,|\Delta_D\Delta_{Q/2}|\,e^{-r/\xi},
\end{equation}
with the pinned phase $\phi_\alpha=-\pi/2$. The crucial factor is
$\cos(\theta-\theta_a)$: it encodes the $d$-wave $2\pi$ winding and forces a
real-space nodal line. The CDW-driven field (paper Eq.~15),
$\rho_\alpha(r)=F_c(r)\cos(Q_\alpha\!\cdot r)$ with $F_c$ peaked at $r=0$ and no
angular factor, has no node. We build both fields on a $256\times256$ real-space
grid ($L=160a$), take $|\mathrm{FFT}|^2$, and test the peak structure at
$q\simeq Q_a/2 = 2\pi/8a$.

\section{Results}
\begin{center}
\begin{tabular}{lcc}
\hline
Signature & PDW-driven & CDW-driven \\
\hline
\# peaks near $Q_a/2$ (on $q_y=0$ cut) & \textbf{2 (split)} & \textbf{1 (single)} \\
peak centers ($q_x/Q_{a/2}$) & $0.9,\ 1.1$ & $1.0$ \\
center dip (center/max) & $0.0$ (zero at $Q/2$) & $1.0$ (peak at $Q/2$) \\
real-space nodal line & yes ($\theta=\pi/2,3\pi/2$) & --- \\
$\mathrm{Re}\,\tilde\nu$ sign change at $Q_a/2$ & yes & yes \\
\hline
\end{tabular}
\end{center}

All four discriminating checks pass. The PDW-driven field shows a clean split
double peak with a \emph{zero exactly at} $q=Q_a/2$ (center/max $=0$), while the
CDW-driven field peaks precisely at $Q_a/2$ (center/max $=1$) --- the qualitative
contrast the paper proposes as the experimental discriminator. The detected
real-space nodal angles ($1.571,\,4.712$ rad) match the predicted
$\theta_a\pm\pi/2$ exactly.

\paragraph{Honest quantitative caveat.} The measured peak splitting on our coarse
grid is $\delta q\approx0.157$ (in units where $Q_a/2=0.785$), whereas the
paper's estimate is $\delta q\sim1/\xi\approx0.067$. The splitting is therefore
reproduced \emph{qualitatively} (double peak, correct symmetry) but is
grid-resolution-limited in magnitude; a finer grid / larger box would sharpen it
toward $1/\xi$. We reproduce the sign and structure of the effect, not a
digit-precise $\delta q$.

\section{Assessment}
The paper's central falsifiable claim --- that a PDW-driven period-8 CDW produces
a \emph{split} Q/2 STM Fourier peak plus a real-space nodal line, distinguishable
from the \emph{single} peak of a CDW-driven scenario --- is reproduced from
scratch (4/4 qualitative checks). We did not reproduce the full BdG
exact-diagonalization LDoS maps of Appendix~B (out of scope for the fast probe),
so this is an honest \textbf{PARTIAL}: the headline discriminator is confirmed,
the microscopic ED and precise $\delta q\sim1/\xi$ magnitude are not.

\section*{Provenance}
Kernel credit: \texttt{loop\_current\_meanfield\_kernel.py} (Ollie), used for
methodological provenance only (real-space order-parameter $\to$ momentum-space
readout). The physics model here is built from scratch for the cuprate PDW/CDW
vortex problem, as the kagome kernel does not apply.
\end{document}
